\chapter{Joint Replacement Surgery}
\index{Joint Replacement Surgery}

\section*{Definition and Overview}
Joint replacement surgery, also called arthroplasty, is an operation to replace a damaged joint with an artificial one, most commonly for advanced osteoarthritis of the hip or knee. The surgeon removes the worn bone surfaces and fixes a prosthesis, made of metal, plastic, or ceramic, to the remaining bone, creating a smooth, pain-free joint. Joint replacement is considered when pain and disability persist despite exercise, weight management, physiotherapy, and medicines. It is a major operation with a short hospital stay, several months of rehabilitation, and excellent long-term results, making it one of the most effective and frequently performed surgeries in many countries.

\section*{Epidemiology}
More than 100,000 hip and knee replacements are performed in many countries each year, and rates are among the highest in the world. The great majority are for osteoarthritis, and the remainder follow rheumatoid arthritis, hip fracture, avascular necrosis, and other conditions. The average age is the late 60s, but younger people increasingly have surgery, and the number of procedures is growing with the ageing population. Shoulder, ankle, and elbow replacements are less common but also increasing. Most implants last 15 to 25 years, after which revision surgery may be needed.

\section*{Aetiology and Pathophysiology}
Arthroplasty is required when the joint surface is destroyed and cannot be repaired. In osteoarthritis, cartilage loss progresses until bone rubs on bone, causing pain, deformity, and stiffness. The artificial joint replaces the damaged surfaces: in a total hip replacement, the femoral head and the socket are replaced; in a total knee replacement, the ends of the thigh and shin bones and the kneecap are resurfaced. The prosthesis is anchored to the bone, either with cement or by bone growing into a porous surface. The goal is to restore a smooth, stable, pain-free joint that moves through a functional range, and modern implants are designed for durability and near-normal function.

\section*{Clinical Features}
\begin{itemize}
    \item Severe joint pain that interferes with walking, sleep, and daily activities
    \item Pain unrelieved by conservative treatment
    \item Stiffness and reduced movement
    \item Difficulty with stairs, chairs, and personal care
    \item Joint deformity and instability
    \item Surgery is usually a total replacement of the hip or knee
    \item The hospital stay is usually two to four days
    \item Crutches or a walker are used for several weeks, with full recovery over months
\end{itemize}

\section*{Differential Diagnosis}
The decision for surgery follows a thorough assessment of alternatives and risks.
\begin{itemize}
    \item \textbf{Conservative care:} exercise, weight loss, physiotherapy, and pain medicines
    \item \textbf{Joint injections:} steroid and lubricant injections
    \item \textbf{Alternative surgery:} arthroscopy, osteotomy, and partial replacement
    \item \textbf{Other causes of pain:} referred pain, bursitis, and vascular disease
    \item \textbf{Medical fitness:} heart, lung, and other conditions affect surgical risk
    \item \textbf{Person's goals and expectations:} surgery is elective and individual
\end{itemize}

\section*{When to Seek Medical Care}
Anyone with severe joint pain that limits life should be assessed for the full range of treatment options. A hot, swollen, painful joint with fever requires urgent care, as does sudden pain or deformity after injury. People with a planned replacement should optimise their health beforehand and seek early review of any new symptoms after surgery.

\textbf{Contact your local emergency services for:}
\begin{itemize}
    \item A hot, swollen, painful joint with fever, which may indicate infection
    \item Sudden chest pain or breathlessness after surgery, which may indicate a blood clot
    \item Severe bleeding from the wound, or a wound that breaks open
    \item Sudden severe pain, deformity, or inability to move the new joint
    \item Collapse, confusion, or severe illness
\end{itemize}

\section*{Diagnosis and Investigations}
\begin{itemize}
    \item \textbf{Clinical assessment:} the severity of pain, functional limitation, and effect on quality of life
    \item \textbf{X-rays:} confirming the degree of joint damage
    \item \textbf{MRI:} in selected cases
    \item \textbf{Preoperative work-up:} blood tests, ECG, and assessment of heart and lung fitness
    \item \textbf{Medicine review:} adjusting blood thinners and other medicines
    \item \textbf{Dental review:} treating any dental infection before surgery
    \item \textbf{Planning:} discussion of implants, anaesthesia, and recovery
\end{itemize}

\section*{Management}
\begin{itemize}
    \item \textbf{Preoperative optimisation:} exercise, weight loss, and stopping smoking
    \item \textbf{Anaesthesia:} general or regional, chosen with the anaesthetist
    \item \textbf{The operation:} removal of damaged surfaces and implantation of the prosthesis
    \item \textbf{Pain management:} medicines and nerve blocks for comfort
    \item \textbf{Early mobilisation:} standing and walking with assistance on the day of or day after surgery
    \item \textbf{Physiotherapy:} a structured program in hospital and at home
    \item \textbf{Prevention of complications:} antibiotics, blood thinners, and wound care
    \item \textbf{Follow-up:} review of the wound, function, and implant at weeks and months
\end{itemize}

\section*{Complications}
\begin{itemize}
    \item Infection of the joint or wound, requiring antibiotics and sometimes further surgery
    \item Blood clots in the leg or lungs
    \item Bleeding and bruising
    \item Dislocation of the hip in the first months
    \item Stiffness and unequal leg lengths
    \item Nerve and blood vessel injury, which is rare
    \item Loosening, wear, or fracture of the implant in the long term
    \item Anaesthesia complications
\end{itemize}

\section*{Prognosis and Outcomes}
Joint replacement surgery has excellent outcomes, with more than 90 per cent of people reporting substantial pain relief and improved function, and high satisfaction rates. Most implants last 15 to 25 years, and revision can be performed when needed. Recovery takes commitment over three to six months, with physiotherapy central to the result. For people with severe arthritis, joint replacement is transformative, restoring mobility, sleep, and independence.

\section*{Prevention and Self-Care}
\begin{itemize}
    \item Lose weight and improve fitness before surgery
    \item Stop smoking and control blood sugar and blood pressure
    \item Prepare the home: raised chairs, clear walkways, and bathroom aids
    \item Follow pre-admission fasting and medicine instructions
    \item Do the prescribed exercises after surgery
    \item Use walking aids until cleared by the team
    \item Keep the wound clean and dry, and report problems
    \item Attend all follow-up appointments
    \item Protect the new joint from falls and high-impact activities as advised
    \item Seek urgent care for fever, chest pain, or a hot, painful joint
\end{itemize}

\section*{Sources and Further Reading}
The following reputable sources provide further information for healthcare students and patients seeking reliable, current advice about joint replacement surgery.
\begin{itemize}
    \item Australian Orthopaedic Association. \textit{Patient information on joint replacement.} https://aoa.org.au/
    \item Healthdirect Australia. \textit{Knee and hip replacement.} https://www.healthdirect.gov.au/
    \item Australian Commission on Safety and Quality in Health Care. \textit{Hip and knee replacement.} https://www.safetyandquality.gov.au/
    \item Arthritis Australia. \textit{Surgery for arthritis.} https://arthritisaustralia.com.au/
    \item NHS. \textit{Hip and knee replacement.} https://www.nhs.uk/
    \item Mayo Clinic. \textit{Joint replacement surgery.} https://www.mayoclinic.org/
\end{itemize}
